\chapter{Grammar Summary}
\label{app:grammar}

This is a condensed form of \file{src/core/javacc/clips.jj}, the JavaCC
grammar the engine is built from. Square brackets mark optional parts,
\fn{*} zero or more repetitions, \fn{+} one or more, and \fn{|}
alternatives. Words in angle brackets are non-terminals; everything else
is literal.

\section{Top level}

\begin{lstlisting}[style=shell]
<input>       ::= <form>*
<form>        ::= ( <expression> ) | ?*global* | <symbol> | <string>
<expression>  ::= <construct> | <special> | <call>
<call>        ::= <function-name> <argument>*
<argument>    ::= <literal> | nil | ?var | ?*global* | $?var | ( <expression> )
                | ( ?var <argument>* )
<literal>     ::= <integer> | <float> | <string> | <symbol>
                | TRUE | true | FALSE | false | nil
\end{lstlisting}

The last form of \fn{<argument>}, a variable and values in parentheses, is
the loop specification of \fn{loop-for-count} and \fn{progn\$}.
The shell, \fn{batch}, \fn{build} and \fn{eval} read \fn{<form>}s.
Data files read by \fn{load-facts}, \fn{load-stream} and \fn{load-graph}
contain \fn{<fact>}s only (below), without \fn{assert}.

\section{Constructs}

\begin{lstlisting}[style=shell]
<construct>   ::= <deftemplate> | <deffacts> | <defrule> | <defquery>
                | <defgraphquery> | <defcube> | <deffunction> | <defmodule>
                | <defclass>

<deftemplate> ::= deftemplate <name> <slot-def>+
<slot-def>    ::= ( slot <name> [ ( type <type> ) ]
                               [ ( default <number> | <string> ) ] )
                | ( multislot <name> )
<type>        ::= INTEGER | SHORT | LONG | FLOAT | DOUBLE | STRING | BOOLEAN
                | symbol | DATE

<deffacts>    ::= deffacts <name> [ <string> ] <fact>+

<defrule>     ::= defrule <name> [ <string> ] [ <declare> ] <ce>* => <action>*
                  [ <<= <action>* ]
<declare>     ::= ( declare [ ( salience <integer> ) ]
                            [ ( auto-focus <bool> ) ]
                            [ ( rule-version <symbol> | <number> ) ]
                            [ ( remember-alpha <bool> ) ]
                            [ ( effective-date <string> ) ]
                            [ ( expiration-date <string> ) ]
                            [ ( chaining-direction <symbol> ) ]
                            [ ( no-agenda <bool> ) ]
                            [ ( hashed-memory <bool> ) ]
                            [ ( temporal-activation <bool> ) ] )
<action>      ::= ( <expression> )

<defquery>    ::= defquery <name> [ <string> ]
                  ( declare ( variables ?var+ ) ) <ce>*
<defgraphquery> ::= defgraphquery <name> [ <string> ]
                  ( declare ( variables ?var+ ) ) <ce>*

<defcube>     ::= defcube <name> ( <- <pattern> )+ ( <dimension> | <measure> )+
<dimension>   ::= ( dimension <name> ?var )
<measure>     ::= ( <name> ( <measure-fn> ?var ) )

<deffunction> ::= deffunction <name> ( <parameter>* ) <body>*
<body>        ::= ( <expression> ) | <literal> | ?var | ?*global*
<parameter>   ::= ?var | ?*global* | $?var
<defmodule>   ::= defmodule <name>
<defclass>    ::= defclass <classname> [ <template-name> [ <parent-template> ] ]
\end{lstlisting}

Rule and template names may be written \fn{MODULE::name}. The properties
in \fn{<declare>} may appear in any order; each is optional.
\fn{defglobal} is parsed as \fn{(defglobal (?*name* [=] [<argument>])*)}.

\section{Conditional elements}

\begin{lstlisting}[style=shell]
<ce>          ::= [ ?var <- ] <pattern>
                | ( not <pattern> )
                | ( not ( and <group-element>+ ) )
                | ( test ( <call> ) )
                | ( and <ce>+ )
                | ( or <ce>+ )
                | ( forall <pattern> <pattern>+ )
                | ( exists <ce>+ )
                | ( only <ce>+ )
                | ( multiple <ce>+ )
                | [ ?var <- ] ( temporal <temporal-declare> <pattern> )
                | [ ?var <- ] ( cubequery ( <cube-name> ( <name> ?var )+ ) )

<temporal-declare> ::= ( declare ?var ( relative-time <number> )
                         [ ( interval-time <number> [ ( <call> ) ] ) ] )

<group-element> ::= <pattern> | ( not <pattern> ) | ( test ( <call> ) )
                  | ( and <group-element>+ ) | ( not ( and <group-element>+ ) )
                  | ( <call> )

<pattern>     ::= ( <template-name> <slot-constraint>* )
<slot-constraint> ::= ( <slot-name> <constraint> )
<constraint>  ::= <literal>
                | ?var
                | $?var
                | ?var&:( <predicate> )
                | =( <call> )
                | ~?var | ~<string> | ~nil
                | <literal> ( | <literal> )+
                | <constraint-atom> ( & <constraint-atom> )+
<constraint-atom> ::= <literal> | ~<literal> | ?var | ~?var
<predicate>   ::= <function-name> <argument>+
\end{lstlisting}

The predicate function name may be a symbol or one of \fn{+ - * / < >
<= >= = == <> != eq and or not}. \fn{\textasciitilde} before a number or an
unquoted symbol is part of the symbol, not a negation. An \fn{or}
element is expanded by the parser into one rule per combination of
alternatives (Section~\ref{sec:or-ce}); \fn{forall} is accepted in
rules only. Every \fn{(not (and (} begins a group
(Section~\ref{sec:not-and}), in which a call is taken as a test; groups
too are accepted in rules only.

\section{Special forms}

These are parsed by dedicated productions rather than as generic calls,
because their arguments are not expressions:

\begin{lstlisting}[style=shell]
<special>     ::= assert <fact>
                | assert-temporal <fact>
                | retract ?var | <integer>
                | modify ?var <slot-value>+
                | duplicate ?var <slot-value>+
                | bind ?var <argument>+ | bind ?*global* <argument>
<fact>        ::= ( <template-name> <slot-value>+ )
<slot-value>  ::= ( <slot-name> <value>* )
<value>       ::= <literal> | ?var | ?*global* | $?var | ( <expression> )
\end{lstlisting}

A slot value is a literal, a variable or a call, evaluated when the fact
is built; a multislot lists several values, or one variable or call that
yields a list. Inside \fn{assert} an empty \fn{(slot)} is accepted and
ignored.

\section{Tokens}

\begin{center}
\begin{tabularx}{\textwidth}{lL}
\toprule
comment & \fn{;;} to end of line \\
integer & optional minus, digits \\
float & optional minus, digits, \fn{.}, digits, optional exponent \\
string & \fn{"..."} or \fn{'...'}; escapes \fn{\textbackslash\textbackslash} and an escaped quote only \\
symbol & letters, digits and \fn{\_ : - \$ . @ /}, plus \fn{*} after the first character; may not start like a number \\
path & a symbol containing backslashes, for Windows file names \\
variable & \fn{?} then a letter or digit, then letters, digits, \fn{-} and \fn{\_} \\
global & \fn{?*name*} \\
multifield & \fn{\$?name} \\
operators & \fn{=> <<= <- + - * / < > <= >= = == <> != eq and or not \& | \textasciitilde} \\
\bottomrule
\end{tabularx}
\end{center}

Reserved words that are tokens but appear in no production, and so are
simply unusable as names: \fn{aggregate}, \fn{allowed-values},
\fn{command}, \fn{describe-class}, \fn{logical}, \fn{no-loop},
\fn{NUMBER}, \fn{send}.
